% GENERATED FILE. Edit canonical lessons, then run npm run generate:books.
% canonical-source-hash: fnv1a64:894289e6b4b77674
% canonical-lessons: AR-C90-huna, AR-C90-hunaka, AR-C90-yamin, AR-C90-matjar, AR-C90-shari

\chapter{Here, There, Right, Shop, Street}
\label{ch:ar-here-there-right-shop-street}
\hlchaptermodality{90}

\begin{chapteropening}
\textbf{By the end of this chapter:} \emph{I can say here, there, right, a shop and a street.}

Complete the last lesson of chapter 90: I can say here, there, right, a shop and a street.
\end{chapteropening}

\section[\texorpdfstring{\ar{هنا} (hunā)}{هنا (hunā)}]{\ar{هنا} (hunā) — here}

\label{lesson:AR-C90-huna}



\begin{quote}
Before the new one: say the Arabic for in front of, then the Arabic for behind.
\end{quote}

\subsection*{You'll want to know: \ar{هنا}}
\textbf{\ar{هنا}} — \emph{hunā} — \textquotedblleft{}here\textquotedblright{}.

\begin{cousinweb}
\textbf{\ar{هنا}} (\emph{hunā}) is an old Arabic pointing word of place. It is not built on a three-letter root; Hebrew \emph{hēnnā}, \textquotedblleft{}hither\textquotedblright{}, is often compared.
\end{cousinweb}

\begin{grammarlens}[title={What you've built}]
One more word for this chapter.
\end{grammarlens}

\subsection*{Your turn}
\begin{itemize}
\raggedright
  \item \emph{Say it:} \emph{hunā}
  \item \emph{Say it:} \emph{hunā}, once more
  \item \emph{Say it:} say \emph{hunā}
  \item \emph{Recall:} say the Arabic for in front of, then the Arabic for behind, then say \emph{hunā} again
\end{itemize}

\begin{tcolorbox}[breakable,colback=teal!4,colframe=teal!35!black,title={Before you move on}]
What does \ar{هنا} mean? (\textquotedblleft{}Here\textquotedblright{}.) Say it once more.
\end{tcolorbox}
\section[\texorpdfstring{\ar{هناك} (hunāka)}{هناك (hunāka)}]{\ar{هناك} (hunāka) — there}

\label{lesson:AR-C90-hunaka}



\begin{quote}
Before the new one: say the Arabic for behind, then the Arabic for here.
\end{quote}

\subsection*{You'll want to know: \ar{هناك}}
\textbf{\ar{هناك}} — \emph{hunāka} — \textquotedblleft{}there\textquotedblright{}.

\begin{cousinweb}
\textbf{\ar{هناك}} (\emph{hunāka}) is built from \emph{hunā}, \textquotedblleft{}here\textquotedblright{}, with the distance ending \emph{-ka}, the same ending heard in \emph{dhālika}, \textquotedblleft{}that\textquotedblright{}.
\end{cousinweb}

\begin{grammarlens}[title={What you've built}]
One more word for this chapter.
\end{grammarlens}

\subsection*{Your turn}
\begin{itemize}
\raggedright
  \item \emph{Say it:} \emph{hunāka}
  \item \emph{Say it:} \emph{hunāka}, once more
  \item \emph{Say it:} say \emph{hunāka}
  \item \emph{Recall:} say the Arabic for behind, then the Arabic for here, then say \emph{hunāka} again
\end{itemize}

\begin{tcolorbox}[breakable,colback=teal!4,colframe=teal!35!black,title={Before you move on}]
What does \ar{هناك} mean? (\textquotedblleft{}There\textquotedblright{}.) Say it once more.
\end{tcolorbox}
\section[\texorpdfstring{\ar{يمين} (yamīn)}{يمين (yamīn)}]{\ar{يمين} (yamīn) — right (side)}

\label{lesson:AR-C90-yamin}



\begin{quote}
Before the new one: say the Arabic for here, then the Arabic for there.
\end{quote}

\subsection*{You'll want to know: \ar{يمين}}
\textbf{\ar{يمين}} — \emph{yamīn} — \textquotedblleft{}right (side)\textquotedblright{}.

\begin{cousinweb}
\textbf{\ar{يمين}} (\emph{yamīn}) is built on the native Arabic root \emph{y-m-n}, \textquotedblleft{}the right hand; good fortune\textquotedblright{} (Lane). Hebrew \emph{yamin}, \textquotedblleft{}right hand\textquotedblright{}, heard in the name Benjamin, is its cousin.
\end{cousinweb}

\begin{grammarlens}[title={What you've built}]
One more word for this chapter.
\end{grammarlens}

\subsection*{Your turn}
\begin{itemize}
\raggedright
  \item \emph{Say it:} \emph{yamīn}
  \item \emph{Say it:} \emph{yamīn}, once more
  \item \emph{Say it:} say \emph{yamīn}
  \item \emph{Recall:} say the Arabic for here, then the Arabic for there, then say \emph{yamīn} again
\end{itemize}

\begin{tcolorbox}[breakable,colback=teal!4,colframe=teal!35!black,title={Before you move on}]
What does \ar{يمين} mean? (\textquotedblleft{}Right (side)\textquotedblright{}.) Say it once more.
\end{tcolorbox}
\section[\texorpdfstring{\ar{متجر} (matjar)}{متجر (matjar)}]{\ar{متجر} (matjar) — a shop}

\label{lesson:AR-C90-matjar}



\begin{quote}
Before the new one: say the Arabic for there, then the Arabic for right.
\end{quote}

\subsection*{You'll want to know: \ar{متجر}}
\textbf{\ar{متجر}} — \emph{matjar} — \textquotedblleft{}a shop\textquotedblright{}.

\begin{cousinweb}
\textbf{\ar{متجر}} (\emph{matjar}) is built on the native Arabic root \emph{t-j-r}, \textquotedblleft{}to trade\textquotedblright{} (Lane). \emph{Matjar} is the \emph{ma-} noun of trading, \textquotedblleft{}a place of trade\textquotedblright{}.
\end{cousinweb}

\begin{grammarlens}[title={What you've built}]
One more word for this chapter.
\end{grammarlens}

\subsection*{Your turn}
\begin{itemize}
\raggedright
  \item \emph{Say it:} \emph{matjar}
  \item \emph{Say it:} \emph{matjar}, once more
  \item \emph{Say it:} say \emph{matjar}
  \item \emph{Recall:} say the Arabic for there, then the Arabic for right, then say \emph{matjar} again
\end{itemize}

\begin{tcolorbox}[breakable,colback=teal!4,colframe=teal!35!black,title={Before you move on}]
What does \ar{متجر} mean? (\textquotedblleft{}A shop\textquotedblright{}.) Say it once more.
\end{tcolorbox}
\section[\texorpdfstring{\ar{شارع} (shāriʿ)}{شارع (shāriʿ)}]{\ar{شارع} (shāriʿ) — a street}

\label{lesson:AR-C90-shari}



\begin{quote}
Before the new one: say the Arabic for right, then the Arabic for a shop.
\end{quote}

\subsection*{You'll want to know: \ar{شارع}}
\textbf{\ar{شارع}} — \emph{shāriʿ} — \textquotedblleft{}a street\textquotedblright{}.

\begin{cousinweb}
\textbf{\ar{شارع}} (\emph{shāriʿ}) is built on the native Arabic root \emph{sh-r-ʿ}, \textquotedblleft{}to enter upon, to open a way\textquotedblright{} (Lane). Lane's \emph{shāriʿ} is an open road; the same root gives \emph{sharīʿa}, \textquotedblleft{}a path, a law\textquotedblright{}.
\end{cousinweb}

\begin{grammarlens}[title={What you've built}]
One more word for this chapter.
\end{grammarlens}

\subsection*{Your turn}
\begin{itemize}
\raggedright
  \item \emph{Say it:} \emph{shāriʿ}
  \item \emph{Say it:} \emph{shāriʿ}, once more
  \item \emph{Say it:} say \emph{shāriʿ}
  \item \emph{Recall:} say the Arabic for right, then the Arabic for a shop, then say \emph{shāriʿ} again
\end{itemize}

\begin{tcolorbox}[breakable,colback=teal!4,colframe=teal!35!black,title={Before you move on}]
What does \ar{شارع} mean? (\textquotedblleft{}A street\textquotedblright{}.) Say it once more.
\end{tcolorbox}
